%==============================================================================
% 7 RESULTS II: CROSS-COUPLING
% Executado em 14/10/2026 (CHANGELOG 2026-10-14 12:11): coupled_mimo_check (primaria,
% atrasos por par, bootstrap, harness) + controles C- (fopdt) e C+ (mujoco-rollers-delay).
% Tabela gerada: tables/tab_journal_coupling (dtssl analysis run journal_tables --kw
% tables=coupling). Fontes: results/coupled_mimo_check_20261014_1211.{md,json},
% results/coupled_mimo_check_control-fopdt_20261014_1211.json,
% results/coupled_mimo_check_control-mujoco-rollers-delay_20261014_1216.json;
% pre-registro docs/PREREGISTRO_MIMO_CHECK_20260923.md (as frases do §8 estao usadas
% no paragrafo M1). Rascunho de 14/10 escrito a partir dos artefatos; revisar a voz.
%==============================================================================
\section{Results II: cross-coupling between channels}
\label{sec:res-coupling}

% ==============================================================================
% tab_journal_coupling: generated by `dtssl analysis run journal_tables` (dtssl.analysis.journal_tables).
% Do not edit by hand: re-run the generator. Sources:
%   results/coupled_mimo_check_20261014_1211.json
%   results/coupled_mimo_check_control-fopdt_20261014_1211.json
%   results/coupled_mimo_check_control-mujoco-rollers-delay_20261014_1216.json
%   results/CHANGELOG.md 2026-10-14 12:11
%   docs/PREREGISTRO_MIMO_CHECK_20260923.md
% ==============================================================================
% Needs \usepackage{booktabs,adjustbox}.
\begin{table*}[t]
\centering
\caption{Cross-coupling check on the competition logs: per-DOF against coupled first-order structure, scored on the frozen hold-out, with the pre-registered verdicts.}
\label{tab:coupling}
\small
\begin{adjustbox}{max width=\textwidth}
\begin{tabular}{lccc}
\toprule
& $v_x$ & $v_y$ & $\omega$ \\
\midrule
pooled hold-out fit (\%): promoted plant & 75.3 & 79.7 & 40.1 \\
\quad diagonal re-fit & 75.5 & 79.9 & 48.8 \\
\quad coupled & 76.5 & 80.1 & 48.9 \\
fit gain, coupled vs. diagonal (p.p.) [95\% CI] & +0.95 [+0.62, +1.07] & +0.23 [+0.02, +0.30] & +0.13 [-0.05, +0.15] \\
cross / direct RMS ratio [95\% CI] & 0.058 [0.046, 0.076] & 0.039 [0.019, 0.059] & 0.048 [0.026, 0.074] \\
direct gain, diagonal / coupled & 1.100 / 1.105 & 1.067 / 1.067 & 1.419 / 1.414 \\
largest cross gain (source) & -0.045 (from $v_y$) & -0.065 (from $\omega$) & -0.047 (from $v_y$) \\
M1 ($\geq 3$ p.p. and ratio $\geq 0.10$): coupling matters & no & no & no \\
\midrule
\multicolumn{4}{l}{free-run $\epos$ at $\hRef$, per-window median (m): promoted 0.0388, diagonal 0.0390, coupled 0.0381} \\
\multicolumn{4}{l}{M3 (gate 8/10, $p<0.05$, $r_b \geq 0.5$): coupled beats diagonal 6/10, $p = 0.46$, $r_b = +0.05$: no} \\
\multicolumn{4}{l}{M2 (half of the excess over 1.0016 open loop): direct yaw gain 1.419 $\rightarrow$ 1.414 with the translational inputs, 1\% removed: falsified} \\
\multicolumn{4}{l}{C$-$ (promoted twin as generator): $G_{\omega\omega}$ 0.9997 / 0.9997 against 1.0017; no spurious coupling: passed} \\
\multicolumn{4}{l}{C$+$ (MuJoCo route b with delay as generator): inflation +7.2\%, below the 10\% premise: not testable} \\
\multicolumn{4}{l}{C$+$, post hoc: $\omega$ from $v_y$ 0.216 of the direct term; $G_{\omega\omega}$ 1.068 $\rightarrow$ 1.075} \\
\bottomrule
\end{tabular}
\end{adjustbox}

\vspace{2pt}
\parbox{\linewidth}{\footnotesize Same estimator (separable least squares on the simulation error), same 36 training events and the same 10 frozen hold-out events for both arms; delays fixed at the promoted plant's 5/4/6 samples. Bootstrap over training events (200 replicates, percentile intervals). The free-run rows use the paper's harness (re-anchored windows, translational bin). The in-silico controls run the same pipeline on noise-free logs generated from the recorded commands.}
\end{table*}


The per-DOF structure of Section~\ref{sec:models} treats $v_x$, $v_y$ and $\omega$ as three independent first-order plants. The natural objection is that it ignores coupling between channels, and the in-silico result of Section~\ref{sec:res-ident} gives that objection a concrete form: in a physical generator, translation induces yaw that is correlated with the yaw command, and the closed-loop yaw gain is inflated by it. We therefore pre-registered three questions before any number was computed on the robot data. M1 (primary): do cross-command terms, added to the same first-order model and fitted with the same estimator on the same 36 training events, improve what the model reproduces on the frozen hold-out? M2 (mechanism): does giving the yaw channel the translational commands as inputs pull the direct yaw gain estimated from the closed-loop logs towards the value of the dedicated open-loop collection? M3 (use): does the coupled model lower the free-run position error of the harness of Section~\ref{sec:protocol}, on the same windows as every other number of this paper? The decision rules, the two in-silico controls and both possible wordings of the outcome were fixed in the pre-registration, and the run followed it without change.

\paragraph{Design.} Both arms share the model $y_i[k] = A_i\,y_i[k-1] + \sum_j B_{ij}\,u_j[k-1-M_{ij}]$ on the 60~Hz grid, the plant alignment of the harness and the delays of the promoted plant (5, 4 and 6 samples for $v_x$, $v_y$ and $\omega$). The diagonal arm sets $B_{ij} = 0$ for $i \neq j$ and is the per-DOF structure re-estimated on the training events; the coupled arm fits the full $3 \times 3$ matrix $B$ with one pole per output. The estimator is the same for both: for a fixed pole the simulated output is linear in the row of $B$, so each row is a separable least-squares problem on the simulation error, the error that the protocol scores, with the pole found by a one-dimensional search. The promoted plant of the paper is scored alongside as the anchor of the published numbers. Coupling is read on two scales at once, the gain in pooled hold-out fit and the ratio between the RMS of the cross terms and the RMS of the direct term, because a cross gain is only interpretable when the term it multiplies contributes to the response. A bootstrap over training events gives the intervals of Table~\ref{tab:coupling}.

\paragraph{M1: coupling does not matter in the envelope of these logs.} A coupled first-order model with cross-command terms, fitted with the same estimator on the same training events, changes the pooled hold-out fit by less than 1~p.p.\ on every DOF, and its cross terms stay below 6\,\% of the direct term's RMS (Table~\ref{tab:coupling}); the largest cross gain, from $v_y$ into $v_x$, is $-0.045$ against a direct gain of 1.105. Neither threshold of the pre-registered rule (3~p.p.\ and a ratio of 0.10 on the same DOF) is reached, and the bootstrap intervals of the fit gains exclude 3~p.p.\ on all three channels. Searching the delay of every input-output pair exhaustively between 0 and 10 samples, the secondary analysis of the pre-registration, moves the pooled fits by at most 1.6~p.p.\ and leaves the verdict unchanged. Within the envelope of these logs the decoupled structure loses nothing measurable; aggressive mixed motion, which the competition logs do not contain, is where a coupled model would have to be tested.

\paragraph{M2: the closed-loop yaw bias is not a linear function of the translational commands.} The premise of the mechanism test is reproduced. Re-estimated on the 36 closed-loop events, the direct yaw gain of the diagonal arm is 1.419 (95\,\% interval 1.360 to 1.517), 41.7\,\% above the 1.0016 of the dedicated open-loop collection and in line with the 1.378 that the prediction-error estimator of Section~\ref{sec:res-ident} gives on the same logs, the correlation bias that direct identification in closed loop is known to carry~\cite{forssell1999closedloop}. Giving the yaw channel the translational commands as inputs leaves the gain at 1.414, which removes 1\,\% of the excess; the rule required half. Restricting the training set to the 13 events of the game label, or searching the delays per pair, does not change the outcome. The hypothesis is falsified in the form it was stated: whatever produces the excess gain is not captured by a term linear in the commanded $v_x$ and $v_y$. The in-silico mechanism of Section~\ref{sec:res-ident} is not thereby refuted, because the pre-registration already stated that a sign-symmetric effect, proportional to $|v|$ for instance, would appear as unexplained variance rather than as a cross gain; the positive control below sharpens that reading.

\paragraph{M3: no gain in use.} In the free-run harness the coupled plant beats the diagonal re-fit in 6 of the 10 hold-out events (one-sided Wilcoxon $p = 0.46$, rank-biserial $+0.05$), short of the gate of 8 of 10 that every other per-event comparison of this paper uses; the per-window medians at $\hRef$ are 0.0388, 0.0390 and 0.0381~m for the promoted, diagonal and coupled plants, within 1~mm of each other. Neither re-fitted arm beats the promoted plant either (5 and 4 of 10 events). The sanity check of the harness reproduces the per-event median of the promoted plant, 0.0421~m, exactly.

\paragraph{Controls.} The negative control runs the same pipeline on noise-free logs generated from the recorded commands by the promoted twin, which has no coupling: both arms recover the yaw gain of the generator to within 0.2\,\% and every cross ratio stays below $10^{-3}$, so the pipeline does not invent coupling and the verdicts on the robot data are interpretable. The positive control replaces the generator by the untuned MuJoCo route~b with the identified delay, the physical generator whose closed-loop yaw inflation Section~\ref{sec:res-ident} reports. With this estimator and all 36 training events the inflation it produces is 7.2\,\% (8.3\,\% on the game label), below the 10\,\% premise fixed in the pre-registration, so the positive control is formally not testable and we report it as such. Its direction is nevertheless informative, and we read it post hoc: the coupled arm does detect the physical coupling, with a cross term from $v_y$ into $\omega$ worth 0.216 of the direct term, the only ratio above 0.10 in the whole study, and still does not lower the direct yaw gain (1.068 to 1.075). A linear cross-command term does not absorb the excess of the direct gain even in a generator where translation-induced yaw is known to exist, and that is the reading we give to the falsification of M2 on the robot.

\paragraph{What the check settles.} The decoupled plant is not a simplification that the data would overturn: at the fidelity the harness measures, the three channels of this robot can be identified separately from operational logs, and the one channel that closed-loop logs identify badly, yaw, is not repaired by coupling but by the short dedicated collection of Section~\ref{sec:res-ident}. The check also bounds what a coupled extension could bring, less than 1~p.p.\ of hold-out fit and no measurable change in free-run error, until data with aggressive mixed motion exist.
